% ---------------------------------------------------------------
% p-adic Variational Quantum Circuits Have No Barren Plateaus
% Target venue: AIQxQIA 2026 (CEUR-WS single-column, single-blind).
% Author: Piero Giacomelli
%
% Compile: pdflatex main && bibtex main && pdflatex main && pdflatex main
% ---------------------------------------------------------------
\documentclass{ceurart}

\sloppy

% ceurart enables microtype font expansion, which needs scalable (Type 1)
% fonts. On systems without the Libertinus fonts the CM fallback is not
% scalable and expansion aborts the build; disabling it is harmless on
% Overleaf and on the CEUR production system.
\microtypesetup{expansion=false}

\usepackage{amsthm}
\usepackage{mathtools}
\usepackage{booktabs}
\usepackage[capitalise,nameinlink]{cleveref}

% ---- theorem environments (ceurart loads amsmath/amssymb but not amsthm) ----
\theoremstyle{plain}
\newtheorem{theorem}{Theorem}[section]
\newtheorem{proposition}[theorem]{Proposition}
\newtheorem{lemma}[theorem]{Lemma}
\newtheorem{corollary}[theorem]{Corollary}

\theoremstyle{definition}
\newtheorem{definition}[theorem]{Definition}
\newtheorem{example}[theorem]{Example}

\theoremstyle{remark}
\newtheorem{remark}[theorem]{Remark}

% Used once, to state the question this paper answers with every term defined
% in the statement itself.
\theoremstyle{plain}
\newtheorem{question}[theorem]{Question}
\crefname{question}{Question}{Questions}

% ---- notation ----
\newcommand{\Qp}{\mathbb{Q}_p}
\newcommand{\Zp}{\mathbb{Z}_p}
\newcommand{\Qpmu}{\mathbb{Q}_{p,\mu}}
\newcommand{\padicnorm}[1]{\left|#1\right|_p}
\newcommand{\opnorm}[1]{\left\|#1\right\|}
\newcommand{\inner}[2]{\langle #1, #2\rangle}
\newcommand{\HH}{\mathcal{H}}
\newcommand{\BB}{\mathcal{B}}
\newcommand{\OK}{\mathcal{O}}
\newcommand{\mm}{\mathfrak{m}}
\newcommand{\Id}{\mathbf{1}}
\newcommand{\loss}{\mathcal{L}}
\newcommand{\ket}[1]{\left|#1\right\rangle}
\newcommand{\bra}[1]{\left\langle #1\right|}
\newcommand{\proj}[1]{\ket{#1}\!\bra{#1}}
\newcommand{\Uni}{\mathrm{U}}
\newcommand{\Sph}{\mathbb{S}}
\newcommand{\FF}{\mathbb{F}}
\newcommand{\Nm}{\mathrm{Nm}}       % field norm z \mapsto z\bar z  (never N)
\DeclareMathOperator{\Tr}{Tr}

\begin{document}

\copyrightyear{2026}
\copyrightclause{Copyright for this paper by its authors.
  Use permitted under Creative Commons License Attribution 4.0
  International (CC BY 4.0).}

\conference{AIQxQIA 2026: International Workshop on AI for Quantum and
  Quantum for AI | co-located with AIxIA 2026, Perugia, Italy}

\title{$p$-adic Variational Quantum Circuits Have No Barren Plateaus}

\author[1]{Piero Giacomelli}[%
email=pgiacome@gmail.com,
]
\address[1]{Tenax SPA, Verona, Italy}

\begin{abstract}
Barren plateaus, the exponential concentration of gradients in random
parameterised quantum circuits, are the central trainability obstruction in
variational quantum machine learning. We show that the phenomenon does not
survive the passage from complex to $p$-adic scalars.

The barren-plateau theorem rests on three ingredients: a compact ensemble
carrying a Haar probability measure, moment data supplied by unitary
$t$-designs, and concentration in the Archimedean absolute value. Over a
quadratic extension $K$ of $\Qp$ each of the three breaks, and they break in
different directions. The first fails outright: the unitary group of $K^{N}$ is
non-compact for $N\ge 2$, so it carries no Haar probability measure and
``Haar-random circuit'' has no meaning; we exhibit an explicit unbounded family
at $p=3$. The compact replacement is the group $\Uni_N(\OK)$ of integral
unitaries, which contains every circuit realisable from the known $p$-adic gate
sets. The second ingredient is not missing but redundant: Haar averages over
$\Uni_N(\OK)$ push forward, level by level, to counting measures on the finite
unitary groups $\Uni_N(\OK/\mm^{k})$, and a character sum over
$\Uni_N(\FF_{p^2})$ supplies in closed form the moment data that the complex
theory obtains from $t$-designs. The third ingredient reverses. For an
$n$-qubit circuit whose observable--generator commutator splits over $\FF_p$,
the gradient attains the largest absolute value the $p$-adic convergence radius
allows, with probability at least $1-p^{-1}-10\,p^{1-r}$, where $r$ measures how
far the observable and the generator are from commuting modulo $p$. The bound
does not depend on $n$, and $r$ grows like $2^{n}$, so the correction is doubly
exponentially small exactly where the complex plateau is most severe. We give a
standard observable--generator pair, Pauli $X$ against a phase rotation, that
achieves $r=2^{n-1}$.

The splitting hypothesis is a real restriction, not a formality: we exhibit a
Hermitian matrix over $\FF_9$ that violates it. Subject to it, trainability in
the $p$-adic model is limited by the conditioning of the optimiser, not by
gradient concentration.
\end{abstract}

\begin{keywords}
  quantum machine learning \sep
  barren plateaus \sep
  variational quantum circuits \sep
  $p$-adic numbers \sep
  non-Archimedean analysis \sep
  trainability
\end{keywords}

\maketitle

% =================================================================
\section{Introduction}
\label{sec:intro}

A variational quantum circuit prepares $U(\theta)\psi_0$ and tunes $\theta$ to
minimise $\loss(\theta)=\inner{U(\theta)\psi_0}{M\,U(\theta)\psi_0}$ for an
observable $M$. McClean et al.~\cite{Mc18} showed that if $U(\theta)$ is drawn
from an ensemble forming a unitary $2$-design on $\mathrm{U}(2^n)$, then
$\mathbb{E}[\partial_\theta\loss]=0$ and
$\mathrm{Var}[\partial_\theta\loss]=O(2^{-n})$. Gradients therefore concentrate
exponentially: the landscape is flat almost everywhere, and the number of
measurements needed to resolve a descent direction grows exponentially in the
qubit count. Later work traced the effect to the cost function~\cite{Ce21} and
to ansatz expressibility~\cite{Ho22}, and it remains the principal argument
against scalable variational quantum machine learning.

Read carefully, the theorem is a statement about the real numbers. It says that
a real quantity, the gradient, is small in the Archimedean absolute value, with
high probability under the Haar measure of a compact group whose volume
concentrates as the dimension grows. Three separate facts are doing the work:
the ensemble is a compact group and so carries a Haar probability measure; the
second moment of that measure is computable because deep circuits approximate
$2$-designs; and the resulting variance is small in an absolute value for which
``small'' is a limiting statement. None of the three is a fact about quantum
mechanics. All three are facts about $\mathbb{C}$.

A parallel line of work replaces those scalars. Following Aniello, Mancini and
Parisi~\cite{AMP23,AMP24,AMP23trace}, one rebuilds quantum kinematics over a
quadratic extension $K=\Qpmu$ of the field $\Qp$ of $p$-adic numbers. States
are vectors in $K^{N}$, gates are unitaries for a $K$-valued Hermitian form,
and the geometry is \emph{ultrametric}:
$\padicnorm{x+y}\le\max\{\padicnorm x,\padicnorm y\}$. The motivation is that
ultrametric geometry is the intrinsic geometry of hierarchical data --
taxonomies, phylogenies, file systems -- which classical $p$-adic learning
theory has exploited for decades~\cite{Kh98,ZuGa22,ZG24dnn,Ma25}. The quantum
side of the programme now has a qubit~\cite{SMW22qub}, a tensor
product~\cite{Mu25tensor}, a universal gate set~\cite{SLMW26gates}, and a first
quantum neural network~\cite{ZGZ26qnn}.

Whether barren plateaus survive the substitution appears to be open, and the
question needs care before it can even be asked. The complex statement concerns
a variance, and a variance decaying like $2^{-n}$ has no $p$-adic meaning: the
$p$-adic absolute value takes only the discrete values $p^{\mathbb{Z}}$, and a
quantity is small when it is divisible by a high power of $p$
(\Cref{sec:prelim}). The faithful translation asks how often the gradient
attains the largest absolute value it can have. We state the question in the
form in which this paper answers it, with every term defined in the statement.

\begin{question}[Barren plateaus over $K$]
\label{qst:plateau}
Let $p$ be an odd prime, let $K$ be the unramified quadratic extension of $\Qp$,
with ring of integers $\OK$ and non-trivial automorphism $z\mapsto\bar z$, and
let $N=2^{n}$. Write $U^\ast$ for the conjugate transpose,
$\inner vw=\sum_i\bar v_iw_i$ for the standard Hermitian form on $K^{N}$, and
let $G_N$ be the group of matrices $U$ with entries in $\OK$ and
$U^\ast U=\Id$, with Haar probability measure $\nu$. Fix an observable
$M=M^\ast$ and a generator $A=-A^\ast$ with entries in $\OK$, a vector
$\psi_0\in\OK^{N}$ with $\inner{\psi_0}{\psi_0}=1$, and a block $U_B\in G_N$.
For $U_A\in G_N$ and $\theta\in\Zp$ consider the circuit and the loss
\[
  U(\theta)\;=\;U_B\,\exp_p(p\,\theta A)\,U_A,
  \qquad
  \loss(\theta)\;=\;\inner{U(\theta)\psi_0}{M\,U(\theta)\psi_0},
\]
where $\exp_p$ is the exponential series, and let $\partial_\theta\loss$ be
the coefficient of $\theta$ in the expansion of $\loss$ at $\theta=0$. Then
$\padicnorm{\partial_\theta\loss}\le p^{-1}$ always. If $U_A$ is drawn from
$\nu$, does
$\Pr\bigl[\,\padicnorm{\partial_\theta\loss}=p^{-1}\,\bigr]$ remain bounded
away from zero as $n\to\infty$, and for which pairs $(M,A)$?
\end{question}

Three features of the statement are forced rather than chosen. The ensemble is
$G_N$ and not the full unitary group of $K^{N}$, because the full group is
non-compact and carries no Haar probability measure, so posed over it the
question has no meaning; this is the first point of the list below. The
generator enters as $p\,\theta A$ because that is exactly what convergence of
$\exp_p$ requires. And the benchmark is $p^{-1}$ rather than $1$, because every
convergent $p$-adic ansatz confines its one-parameter subgroups to the
principal congruence subgroup, so $p^{-1}$ is the largest gradient any such
ansatz permits (\Cref{rem:scale}).

A heuristic suggests the answer is yes. The normalised gradient
$p^{-1}\partial_\theta\loss$ lies in $\Zp$, and a Haar-random element of $\Zp$
is a unit with probability $1-p^{-1}$, so one expects
$\padicnorm{\partial_\theta\loss}=p^{-1}$ with probability close to $1-p^{-1}$,
whatever the qubit count. \Cref{thm:main} confirms this, uniformly in $n$,
under a splitting hypothesis on the commutator $[U_B^\ast MU_B,A]$ reduced
modulo $p$ (\Cref{def:nondeg}). The configurations that fail are
those in which $U_B^\ast MU_B$ and $A$ commute modulo $p$ (\Cref{cor:deg}); as
$U_B$ varies they form a clopen union of congruence classes of measure
$O(p^{-1})$ (\Cref{sec:discussion}), small but strictly positive, a
distinction the complex theory never has to make. Outside the splitting class
the question remains open (\Cref{sec:limitations}).

We take the three ingredients in turn.

\begin{itemize}
\item \textbf{The ensemble is wrong} (\Cref{sec:noncompact}). $\Uni(K^{N})$ is
  non-compact for every $N\ge2$; we give an explicit unbounded family at $p=3$
  (\Cref{prop:noncompact,ex:p3}). There is therefore no Haar probability
  measure on the full unitary group, and the plateau question, phrased over it,
  is ill posed. The compact replacement is the integral unitary group
  $G_N=\Uni_N(\OK)$, which contains every circuit built from the currently
  known realisable $p$-adic gate sets (\Cref{prop:gates}).
\item \textbf{The missing combinatorics are not needed}
  (\Cref{sec:counting}). Haar averages over $G_N$ push forward to uniform
  counting measures on the finite unitary groups $\Uni_N(\OK/\mm^{k})$
  (\Cref{lem:pushforward}), and Haar-random states push forward to the uniform
  measure on a finite Hermitian sphere (\Cref{lem:transitive}). At level one
  this turns every question about leading $p$-adic digits into a character sum
  over $\FF_{p^2}$, which we evaluate in closed form (\Cref{lem:count}).
\item \textbf{The conclusion reverses} (\Cref{sec:main}).
  \Cref{thm:main} shows that for an $n$-qubit circuit whose
  observable--generator commutator is $r$-non-degenerate modulo $p$ in the sense
  of \Cref{def:nondeg},
  \[
    \Pr\bigl[\,\padicnorm{\partial_\theta\loss}=p^{-1}\,\bigr]\;\ge\;
    1-\tfrac1p-10\,p^{1-r},
  \]
  and $p^{-1}$ is the largest absolute value any convergent $p$-adic ansatz
  permits (\Cref{rem:scale}). The gradient is therefore generically as large as
  it can be, with a bound that does not depend on $n$ at all.
\end{itemize}

\Cref{tab:compare} summarises the comparison. Two features of the $p$-adic
answer are worth flagging early because they have no complex analogue. The
bound improves as the circuit becomes more expressive, inverting the
relationship established in~\cite{Ho22}; and the set of bad
observable--generator pairs is not a vague ``exponentially large region'' but
the explicit locus $[\,\widehat M,A\,]\equiv0\bmod p$, a condition one can check
by reducing two integer matrices and multiplying them.

\begin{table}[t]
\caption{Barren plateaus over $\mathbb{C}$ and over $K=\Qpmu$. Each row is one
  ingredient of the argument of~\cite{Mc18} and what becomes of it.}
\label{tab:compare}
\centering
\small
\begin{tabular}{@{}p{0.26\linewidth}p{0.32\linewidth}p{0.32\linewidth}@{}}
\toprule
 & \textbf{Complex} & \textbf{$p$-adic} \\
\midrule
Ensemble & $\mathrm{U}(2^{n})$, compact & $\Uni(\HH_N)$ non-compact; use
  $G_N=\Uni_N(\OK)$ \\
Source of moments & unitary $t$-designs & character sums over
  $\Uni_N(\FF_{p^2})$ \\
Gradient statistic & variance $O(2^{-n})$ & $\Pr[\text{maximal}]\ge
  1-p^{-1}-10p^{1-r}$ \\
Scaling in $n$ & exponential decay & none \\
Bad set & full measure as $n\to\infty$ & the locus
  $[\widehat M,A]\equiv0\bmod p$ \\
Role of expressibility & more expressive $\Rightarrow$ flatter~\cite{Ho22} &
  more random $\Rightarrow$ larger gradients \\
\bottomrule
\end{tabular}
\end{table}
 We fix $p$ odd and $K/\Qp$ unramified throughout, work in
finite dimension, and randomise a single parameter in the standard two-block
ansatz of~\cite{Mc18}. \Cref{sec:discussion} states precisely what is and is
not established. The load-bearing restriction is the splitting hypothesis of
\Cref{def:nondeg}, which \Cref{rem:notspectral} shows is not automatic; the
example of \Cref{ex:witness} shows it is nevertheless satisfied by the most
ordinary pair one might write down. The paper is theoretical: no $p$-adic
quantum hardware exists, and, as recalled in \Cref{rem:nosampling}, the
$p$-adic model has no Monte Carlo readout, so the empirical protocol behind
barren-plateau plots in the complex setting has no counterpart here.
\Cref{sec:why} argues that the result is nonetheless useful, and that its
first use is diagnostic and concerns the complex theory: the plateau survives
every change we make to the circuit and not the one we make to the field, which
locates the hypothesis a mitigation has to attack.

% =================================================================
\section{Preliminaries}
\label{sec:prelim}

We compress the background to what the proofs use; \cite{Go97} is the reference
for $p$-adic arithmetic and \cite{AMP23,AMP24} for the quantum formalism. A
reader who works on variational algorithms but not on non-Archimedean analysis
needs three facts, and we state them informally before the definitions.

\emph{Small means divisible.} The $p$-adic absolute value measures divisibility,
not magnitude: $\padicnorm x=p^{-v_p(x)}$, where $v_p(x)$ counts the factors of
$p$ in $x$. Taking $p=3$, we have $\left|18\right|_3=1/9$ because $18=2\cdot3^2$,
while $\left|1/9\right|_3=9$ and $\left|2\right|_3=1$. Integers divisible by a high power of
$p$ are the small ones.

\emph{The scale is discrete.} The values taken by $\padicnorm\cdot$ are the
powers $p^{\mathbb{Z}}$, so there is no continuum of small quantities and no
notion of a quantity that is merely somewhat small. A gradient is either a unit,
or divisible by $p$, or by $p^2$, with nothing in between. This is why
``exponential concentration'' has to be re-read before it can be asked about:
$\mathrm{Var}=O(2^{-n})$ has no $p$-adic meaning, and the faithful translation
is whether the valuation $v_p$ of the gradient grows with $n$.
\Cref{qst:plateau} asks the sharpest form of this, which is why the answer here
is a probability of attaining one specific value rather than a decay rate.

\emph{Sums do not grow, and cancellation is a congruence.} The ultrametric
inequality $\padicnorm{x+y}\le\max\{\padicnorm x,\padicnorm y\}$ holds with
equality whenever the two terms have different sizes. A sum is therefore as
large as its largest term unless the leading digits cancel exactly, and that
cancellation is a condition modulo $p$ -- a finite, checkable condition rather
than a limiting one. This is the structural reason the whole second-moment
analysis of \Cref{sec:counting} collapses onto counting in finite groups, where
the complex theory needs $t$-design combinatorics: reduction modulo $p$ carries
Haar measure on the integral unitary group to the uniform measure on a finite
unitary group, and questions about leading digits become character sums over
$\FF_{p^2}$.

One expectation does not carry over, and it is a basic one. Generalised
probabilities in this formalism lie in $\Qp$ rather than in $[0,1]$, so
amplitudes have no frequentist reading and
cannot be estimated by repeated measurement; readout is algebraic, and
\Cref{rem:nosampling} states what a plateau means once sampling is unavailable.

\subsection{Arithmetic}

Fix an odd prime $p$. For $x\in\mathbb{Q}^\ast$ let $v_p(x)$ be the exponent of
$p$ in $x$ and $\padicnorm x=p^{-v_p(x)}$, with $\padicnorm 0=0$ and
$v_p(0)=+\infty$; $\Qp$ is the completion of $\mathbb{Q}$ for this absolute
value and $\Zp=\{\padicnorm x\le1\}$ its ring of integers. The absolute value
is \emph{ultrametric}: $\padicnorm{x+y}\le\max\{\padicnorm x,\padicnorm y\}$,
with equality when $\padicnorm x\ne\padicnorm y$. Its value group is $p^{\mathbb{Z}}$,
a discrete subgroup of $\mathbb{R}_{>0}$; this discreteness is used repeatedly
below and is the single most important difference from the real case.

Let $\mu\in\Zp^\times$ be a non-square unit and $K=\Qpmu=\Qp(\sqrt\mu)$, the
unramified quadratic extension, with involution
$\sigma(x+\sqrt\mu\,y)=x-\sqrt\mu\,y$, written $\bar z$. The absolute value
extends by $\padicnorm z=\padicnorm{\Nm(z)}^{1/2}$, where
\[
  \Nm:K\longrightarrow\Qp,\qquad \Nm(z)=z\bar z
\]
is the field norm. Write $\OK$ for the ring of integers of $K$ and $\mm=p\OK$
for its maximal ideal, so that
\[
  \OK/\mm\;=\;\FF_{p^2},\qquad \Zp/p\Zp\;=\;\FF_p ,
\]
$\sigma$ reduces to the Frobenius $x\mapsto x^{p}$ of $\FF_{p^2}/\FF_p$, and
$\Nm$ reduces to the norm $z\mapsto z^{1+p}$ of $\FF_{p^2}/\FF_p$. For an
integral $z$ we write $\widetilde z\in\FF_{p^2}$ for its residue, and we apply
the same notation entrywise to vectors and matrices. Because $K/\Qp$ is
unramified, $\Nm(K^\ast)$ is exactly the set of elements of even
valuation~\cite{Se73}.

We isolate one arithmetic fact, used three times.

\begin{lemma}[Unit norms]
\label{lem:unitnorm}
$\Nm(\OK^\times)=\Zp^\times$.
\end{lemma}

\begin{proof}
The inclusion $\subseteq$ is clear. Conversely let $c\in\Zp^\times$. The norm
of $\FF_{p^2}/\FF_p$ is the map $z\mapsto z^{1+p}$ on $\FF_{p^2}^\times$; its
kernel has order $p+1$, so its image has order $(p^2-1)/(p+1)=p-1$ and is all
of $\FF_p^\times$. Pick $s_1\in\OK^\times$ with $\Nm(s_1)\equiv c\bmod p$. Then
$c/\Nm(s_1)\in1+p\Zp$, and every element of $1+p\Zp$ is a square in $\Zp$ for
$p$ odd, by Hensel's lemma applied to $X^2-c/\Nm(s_1)$ at $X=1$. Writing
$c/\Nm(s_1)=X^2$ with $X\in1+p\Zp$ and setting $s=Xs_1$ gives
$\Nm(s)=X^2\Nm(s_1)=c$.
\end{proof}

\subsection{States, gates, gradients}

Let $\HH_N=K^{N}$ with $\opnorm{v}=\max_i\padicnorm{v_i}$ and
$\inner vw=\sum_i\bar v_i w_i$, linear in the second argument. This is the
model finite-dimensional $p$-adic Hilbert space~\cite{AMP23}; the $n$-qubit
space is $\HH_N$ with $N=2^n$, built from $\HH_2^{\otimes n}$ via the $p$-adic
tensor product~\cite{Mu25tensor}. An operator $A\in\BB(\HH_N)$ has
$\opnorm A=\max_{ij}\padicnorm{A_{ij}}$ and adjoint
$(A^\ast)_{ij}=\overline{A_{ji}}$; ultrametricity makes $\opnorm{\cdot}$
submultiplicative, since
$\padicnorm{(AB)_{ij}}=\padicnorm{\sum_kA_{ik}B_{kj}}
\le\max_k\padicnorm{A_{ik}}\padicnorm{B_{kj}}$. A \emph{gate} is a $U$ with
$U^\ast U=\Id$; these form the group $\Uni(\HH_N)$. A vector $\psi$ is
\emph{admissible} if $\inner\psi\psi=1$, and observables are selfadjoint
operators.

Two features of the model matter below. First, positivity is unavailable:
$\inner\psi\psi\ge0$ is meaningless because $K$ is unordered, so states form an
affine rather than a convex set and measurement is described by selfadjoint
operator-valued measures with $\Qp$-valued generalised
probabilities~\cite{AMP23,Kh97}. Second:

\begin{remark}[No Monte Carlo readout~\cite{AMP23}]
\label{rem:nosampling}
Generalised probabilities lie in $\Qp$, not $[0,1]$, and cannot be estimated by
repeated sampling. Readout is algebraic: the algorithm returns $p$-adic numbers
which the user inspects digit by digit. The operational meaning of a barren
plateau changes accordingly. In the complex setting a small gradient is bad
because it takes many shots to resolve; in the $p$-adic setting there are no
shots, and a gradient of high valuation is bad because the leading digits of
the parameter update vanish.
\end{remark}

For the one-parameter dependence we use the exponential of the $p$-adic Lie
algebra. The series $\exp_p(X)=\sum_{m\ge0}X^m/m!$ converges whenever
$\opnorm X<p^{-1/(p-1)}$~\cite{DDMS99}, this being the standard bound we adopt,
and for $A^\ast=-A$ it defines a unitary. Following~\cite{Mc18} we take the two-block ansatz, in the
normalisation forced by \Cref{rem:scale} below:
\begin{equation}
\label{eq:ansatz}
  U(\theta)\;=\;U_B\,\exp_p(p\,\theta A)\,U_A ,\qquad
  \theta\in\Zp,\quad A^\ast=-A,\quad A\in \BB(\HH_N)\ \text{integral} .
\end{equation}
Since $\loss$ is $p$-adic analytic in $\theta$, we define $\partial_\theta\loss$
as the coefficient of $\theta$ in its Taylor expansion at $\theta=0$.

\begin{remark}[The scale $p$ is forced, not chosen]
\label{rem:scale}
For $p$ odd we have $0<\tfrac1{p-1}<1$, and the value group of
$\opnorm{\cdot}$ is the \emph{discrete} set $p^{\mathbb{Z}}$. No operator norm
lies strictly between $p^{-1}$ and $1$, so the convergence bound
$\opnorm{X}<p^{-1/(p-1)}$ is not a mild constraint that could be relaxed by a
better estimate: it is \emph{equivalent} to $\opnorm X\le p^{-1}$, that is, to
$X\equiv0\bmod\mm$. Under it, every term $X^{m}/m!$ with $m\ge1$ has norm at
most $p^{-m}\,p^{(m-1)/(p-1)}\le p^{-1}$, since $v_p(m!)\le(m-1)/(p-1)$, so
\[
  \exp_p(\theta X)\;\equiv\;\Id \pmod{\mm}\qquad\text{for all }\theta\in\Zp .
\]
The one-parameter subgroup thus lies inside the principal congruence subgroup
$\ker(\pi_1)$ of \Cref{sec:counting}, and every Taylor coefficient of $\loss$
of positive order has valuation at least $1$. In particular
\[
  \padicnorm{\partial_\theta\loss}\;\le\;p^{-1}
\]
for any integral observable and any ansatz meeting the convergence bound. The generator in~\eqref{eq:ansatz} is written $pA$ with $A$
integral for exactly this reason: the factor $p$ is not a normalisation we are
free to choose, it is the entire content of the convergence condition, and
$p^{-1}$ rather than $1$ is the correct benchmark for ``as large as possible''.
We write
\begin{equation}
\label{eq:normgrad}
  \nabla\loss\;:=\;p^{-1}\,\partial_\theta\loss
\end{equation}
for the \emph{normalised gradient}, which does satisfy
$\padicnorm{\nabla\loss}\le1$, and state our results in terms of it.
\end{remark}

% =================================================================
\section{The ensemble: the unitary group is not compact}
\label{sec:noncompact}

The barren-plateau question presupposes a Haar \emph{probability} measure on
the group of $n$-qubit unitaries. Over $\mathbb{C}$ this is automatic because
$\mathrm{U}(2^n)$ is compact. Over $K$ it fails, and it fails badly.

\begin{proposition}
\label{prop:noncompact}
For every $N\ge 2$ the group $\Uni(\HH_N)$ is unbounded in the operator norm,
hence non-compact, and carries no Haar probability measure.
\end{proposition}

\begin{proof}
It suffices to treat $N=2$ and pad with the identity, since
$U\mapsto U\oplus\Id_{N-2}$ embeds $\Uni(\HH_2)$ isometrically in
$\Uni(\HH_N)$. We construct, for each $m\ge1$, a unitary of operator norm
$p^{m}$.

Put $a=p^{-m}$ and $t=\Nm(a)=p^{-2m}$. Since $\padicnorm t=p^{2m}>1=\padicnorm
1$, ultrametricity gives $\padicnorm{1-t}=\padicnorm t$, so $v_p(1-t)=-2m$ is
even. Because $K/\Qp$ is unramified, $\Nm(K^\ast)$ consists exactly of the
elements of even valuation, so there is $c\in K$ with $\Nm(c)=1-t$. The vector
$u=(a,c)^{\mathsf T}$ then satisfies
\[
  \inner uu \;=\; \bar aa+\bar cc\;=\;\Nm(a)+\Nm(c)\;=\;t+(1-t)\;=\;1 ,
\]
so $u$ is admissible, while $\padicnorm a=p^{m}$.

Both $e_1$ and $u$ are admissible, so $e_1\mapsto u$ is an isometry between
one-dimensional subspaces of the non-degenerate Hermitian space $\HH_2$. Witt's
extension theorem for Hermitian forms over a field with involution, in
characteristic $\ne2$~\cite[Ch.~7]{Gr02}, extends it to an isometry $U$ of
$\HH_2$, that is, to an element of $\Uni(\HH_2)$. Its first column is $u$, so
$\opnorm U\ge\padicnorm a=p^{m}$.

Finally, $\Uni(\HH_N)$ is a closed subgroup of $\mathrm{GL}_N(K)$, hence
locally compact; the compact subsets of $K^{N\times N}$ are the closed bounded
ones, so an unbounded closed subgroup is not compact; and a locally compact
group admits a finite Haar measure only when it is compact.
\end{proof}

\begin{example}[$p=3$]
\label{ex:p3}
Take $\mu=-1$, a non-square unit modulo $3$, so that $K=\mathbb{Q}_3(i)$ is
unramified and $\Nm(x+iy)=x^2+y^2$. With $m=1$ we need $\Nm(a)=1/9$ and
$\Nm(c)=8/9$. Take $a=1/3$ and $c=\tfrac13(x_1+iy_1)$ with $x_1^2+y_1^2=8$: the
congruence $x_1^2+y_1^2\equiv2\pmod3$ has the solution $x_1\equiv y_1\equiv1$,
which lifts to $\Zp$ by Hensel's lemma applied to $X^2-7$ at $X=1$, whose
derivative $2X$ is a unit there. Then $\padicnorm a=\padicnorm c=3$, and
completing $(a,c)^{\mathsf T}$ to a unitary gives an element of $\Uni(\HH_2)$
of operator norm $3$. Replacing $1/3$ by $3^{-m}$ gives operator norm $3^{m}$.
\end{example}

\Cref{prop:noncompact} is not an artefact of a poor choice of Hermitian form.
Over a local field the standard form $\sum\bar z_iz_i$ is isotropic in every
dimension $N\ge2$ -- for $N=2$ isotropy is exactly the statement
$-1\in \Nm(K^\ast)$, which holds in the unramified case -- and unitary groups
of isotropic Hermitian forms are non-compact~\cite[Ch.~3]{PR94}. The
compactness that the complex theory takes for granted comes from the
positive-definiteness of $\sum\bar z_iz_i$ over $\mathbb{C}$, and positivity is
precisely what the $p$-adic model gives up.

\begin{definition}[Integral unitary group]
\label{def:integral}
$G_N \;=\; \Uni_N(\OK)\;=\;\{U\in\Uni(\HH_N)\;:\;U_{ij}\in\OK\ \text{for all
$i,j$}\}$.
\end{definition}

\begin{proposition}
\label{prop:compact}
$G_N$ is a compact group and carries a unique Haar probability measure $\nu$.
\end{proposition}

\begin{proof}
$G_N$ is closed under products, and under inverses because
$U^{-1}=U^\ast$ has entries $\overline{U_{ji}}\in\OK$. As a subset of
$\OK^{N\times N}$ it is bounded, and it is closed, being the solution set of
the polynomial system $U^\ast U=\Id$ over the compact ring $\OK$. A closed
bounded subset of $\OK^{N^2}$ is compact, and a compact group has a unique
normalised Haar measure.
\end{proof}

Restricting to $G_N$ costs nothing for quantum computing, because everything
one can currently build already lies there.

\begin{proposition}
\label{prop:gates}
Every gate whose matrix entries lie in $\OK$ generates, under composition and
tensor product, a subgroup of $G_N$. This applies to the $p$-adic Pauli gates
and CNOT of~\cite{AMP23}, to the gate set drawn from irreducible
representations of $\mathrm{SO}(3)_p$ of~\cite{SLMW26gates}, and  when
$\sqrt2\in K$  to the $p$-adic Hadamard gate.
\end{proposition}

\begin{proof}
Integrality is preserved by products and tensor products, and $G_N$ is a group
by \Cref{prop:compact}. For the three gate sets it remains only to observe that
their entries are integral. The Pauli and CNOT entries lie in $\{0,\pm1\}$. The
entries of the $\mathrm{SO}(3)_p$ gates are roots of unity in $K$ and zeros;
roots of unity are units, hence integral. For the Hadamard gate, if
$\sqrt2\in K$ then $\padicnorm{\sqrt 2}=\padicnorm 2^{1/2}=1$ for odd $p$, so
$1/\sqrt2\in\OK^\times$ and all four entries are units.
\end{proof}

\begin{remark}
\label{rem:ensemble}
The contrast with $\mathbb{C}$ runs deeper than a change of group. There,
``Haar-random unitary'' and ``deep random circuit'' converge to each other, and
that convergence is what makes the $t$-design argument of~\cite{Mc18} bite.
Over $K$ the two notions come apart at the outset: the group of all unitaries
is too large to randomise over, and the group generated by realisable gates is
a compact subgroup with a different measure-theoretic structure altogether. We
take $G_N$ as the ensemble because it is the smallest natural compact group
containing all realisable circuits. Sharper results should be available for the
compact group actually generated by a fixed gate set, which is where a
hardware-aware analysis would have to go.
\end{remark}

% =================================================================
\section{Haar averages are finite counting problems}
\label{sec:counting}

The argument of~\cite{Mc18} rests on unitary $t$-design combinatorics, for
which we know of no $p$-adic counterpart. They are not needed. Compactness of
$G_N$ arrives together with a filtration that turns every question about the
leading digits of a Haar average into a finite count.

For $k\ge1$ write $\OK_k=\OK/\mm^{k}$ and let
$\pi_k:G_N\to \Uni_N(\OK_k)$ be reduction of matrix entries modulo $\mm^k$,
where $\Uni_N(\OK_k)$ is the unitary group of the reduced Hermitian form over
the finite ring $\OK_k$. For $k=1$ this is the finite unitary group
$\Uni_N(\FF_{p^2})$.

\begin{lemma}[Pushforward]
\label{lem:pushforward}
For every $k\ge1$ the map $\pi_k$ is a surjective group homomorphism and
$(\pi_k)_\ast\nu$ is the uniform probability measure on $\Uni_N(\OK_k)$.
Consequently, if $F:G_N\to\Zp$ is a function whose residue $F\bmod p^{k}$
depends on $U$ only through $\pi_k(U)$, then
\[
  \nu\bigl(\{U: v_p(F(U))\ge k\}\bigr)
  \;=\;\frac{\#\{g\in \Uni_N(\OK_k)\;:\;F(g)\equiv0\}}
            {\bigl|\Uni_N(\OK_k)\bigr|}.
\]
\end{lemma}

\begin{proof}
Reduction modulo $\mm^k$ is a ring homomorphism $\OK\to\OK_k$ commuting with
$\sigma$, because $\sigma(\mm^k)=\mm^k$; it therefore carries the relation
$U^\ast U=\Id$ to its reduction, and $\pi_k$ is a homomorphism.

For surjectivity, take $g\in\Uni_N(\OK_k)$ and lift its entries arbitrarily to
a matrix $g_0\in\BB(\HH_N)$ with integral entries. Then $g_0^\ast g_0$ is
selfadjoint and congruent to $\Id$ modulo $\mm^{k}$, say
\[
  g_0^\ast g_0\;=\;\Id+p^{\,l}S,\qquad l=k\ge1,\quad S^\ast=S\ \text{integral} .
\]
We improve $l$ by Newton iteration. Set $T=-S/2$, which is integral because $p$
is odd, and selfadjoint because $S$ is; put $g_1=g_0(\Id+p^{\,l}T)$. Expanding,
and using $T^\ast=T$ in the leftmost factor,
\[
  g_1^\ast g_1
  =(\Id+p^{l}T)(\Id+p^{l}S)(\Id+p^{l}T)
  =\Id+p^{l}\bigl(T+S+T\bigr)+p^{2l}\bigl(TS+ST+T^2\bigr)+p^{3l}TST ,
\]
and $T+T=-S$ kills the order-$p^{l}$ term, so $g_1^\ast g_1=\Id+p^{2l}S'$ with
$S'$ selfadjoint and integral. Moreover $g_1\equiv g_0\bmod \mm^{l}$. Iterating
produces a sequence $g_0,g_1,g_2,\dots$ with $g_{j+1}\equiv g_j\bmod
\mm^{\,2^{j}k}$, which is Cauchy in the complete space $\BB(\HH_N)$ of integral
matrices; its limit $h$ is integral, satisfies $h^\ast h=\Id$  hence is
invertible with $h^{-1}=h^\ast$ integral, so $h\in G_N$  and satisfies
$h\equiv g_0\bmod\mm^{k}$, that is $\pi_k(h)=g$.

Being a continuous surjective homomorphism from the compact group $G_N$ onto
the finite group $\Uni_N(\OK_k)$, $\pi_k$ carries $\nu$ to a translation-%
invariant probability measure on the target, which by uniqueness of Haar
measure is the uniform one. The displayed identity is the definition of the
pushforward applied to the indicator of $\{F\equiv0\bmod p^k\}$, which by
hypothesis is a $\pi_k$-pullback.
\end{proof}

We shall also need that Haar measure on $G_N$ moves a fixed admissible vector
to the uniform measure on the integral sphere
\[
  \Sph_N \;=\;\{u\in\OK^{N}\;:\;\inner uu=1\},\qquad
  \widetilde\Sph_N \;=\;\{\tilde u\in\FF_{p^2}^{N}:\inner{\tilde u}{\tilde
  u}=1\} .
\]
The proof below is elementary and self-contained; it replaces an appeal to
smoothness of the unitary group scheme, and it makes visible exactly where $p$
odd is used.

\begin{lemma}[Transitivity]
\label{lem:transitive}
Let $N\ge1$. Then $G_N$ acts transitively on $\Sph_N$, and for
$\psi_0\in\Sph_N$ the pushforward of $\nu$ under $U\mapsto U\psi_0$ is the
unique $G_N$-invariant probability measure on $\Sph_N$. Reduction modulo $\mm$
maps $\Sph_N$ onto $\widetilde\Sph_N$ and carries that measure to the uniform
measure on $\widetilde\Sph_N$.
\end{lemma}

\begin{proof}
We argue in four steps.

\emph{(i) Every $u\in\Sph_N$ is unimodular.} From
$1=\padicnorm{\inner uu}\le\max_i\padicnorm{u_i}^2\le1$ we get
$\max_i\padicnorm{u_i}=1$.

\emph{(ii) A norm-one vector splits off an orthogonal complement.} Let $L$ be
any free Hermitian $\OK$-lattice whose Gram determinant is a unit  call such
a lattice \emph{unimodular}  and let $u\in L$ with $\inner uu=1$. The map
$P(x)=x-\inner ux\,u$ is $\OK$-linear, because $\inner u\cdot$ is linear in its
second argument, and $\inner u{P(x)}=\inner ux(1-\inner uu)=0$. Hence
$L=\OK u\oplus u^{\perp}$ orthogonally, where $u^{\perp}=\ker\inner u\cdot$. As
a direct summand of a free module over the local ring $\OK$, the lattice
$u^{\perp}$ is free, of rank one less. It is again unimodular: the change of
basis from a basis of $L$ to one adapted to the splitting multiplies the Gram
determinant by $\Nm(\det Q)$ for some $Q\in\mathrm{GL}(\OK)$, hence by a unit,
while the adapted Gram matrix is $\mathrm{diag}(1,G)$, so $\det
G\in\OK^\times$. Applied to $L=\OK^{N}$ with the standard form, whose Gram
determinant is $1$, this splits $\OK^{N}=\OK u\perp u^{\perp}$ with
$u^{\perp}$ unimodular of rank $N-1$.

\emph{(iii) Every unimodular Hermitian $\OK$-lattice $L$ of rank $\ge1$
contains a vector of norm $1$.} The reduction $L/\mm L$ carries a
non-degenerate Hermitian form over $\FF_{p^2}$. That form cannot vanish on the
diagonal: if $\inner xx=0$ for all $x$, then polarising gives
$\Tr_{\FF_{p^2}/\FF_p}\inner xy=0$ for all $x,y$, and replacing $y$ by
$\sqrt\mu\,y$ gives $\Tr_{\FF_{p^2}/\FF_p}(\sqrt\mu\inner xy)=0$; writing
$\inner xy=a+b\sqrt\mu$ these two conditions read $2a=0$ and $2b\mu=0$, so
$\inner xy=0$ identically, contradicting non-degeneracy. Choose therefore
$x_0\in L$ whose residue has non-zero norm, so that
$c=\inner{x_0}{x_0}\in\Zp^\times$. By \Cref{lem:unitnorm} there is
$s\in\OK^\times$ with $\Nm(s)=c^{-1}$, and then $\inner{sx_0}{sx_0}=\Nm(s)c=1$.

\emph{(iv) Transitivity.} Any two unimodular Hermitian $\OK$-lattices of equal
rank are isometric, by induction on the rank: each contains a vector of norm
$1$ by (iii), and (ii) splits that vector off leaving unimodular lattices of
rank one less, the rank-zero case being trivial. Given $u,u'\in\Sph_N$, apply
(ii) to both and choose an isometry
$\phi:u^{\perp}\to u'^{\perp}$; the map sending $u\mapsto u'$ and acting as
$\phi$ on $u^\perp$ is an isometry of $\OK^N$, that is an element $U\in G_N$
with $Uu=u'$.

Uniqueness of the invariant probability measure on a transitive compact group
action is standard. For the reduction, note first that every
$\tilde u\in\widetilde\Sph_N$ lifts: choose any integral $u_0$ with residue
$\tilde u$, so $c=\inner{u_0}{u_0}\in1+p\Zp$; by Hensel $c^{-1}=X^2$ with
$X\in1+p\Zp$, and $u=Xu_0$ lies in $\Sph_N$ with the same residue. Reduction
therefore maps $\Sph_N$ onto $\widetilde\Sph_N$, and it intertwines the action
of $G_N$ with the action of $\Uni_N(\FF_{p^2})$, which is surjective by
\Cref{lem:pushforward}. The image of an invariant measure under an equivariant
surjection onto a transitive finite action is the uniform measure.
\end{proof}

% =================================================================
\section{Gradient valuations}
\label{sec:main}

\subsection{The gradient is a Hermitian form}

\begin{lemma}
\label{lem:grad}
For the ansatz~\eqref{eq:ansatz} with $\psi=U_A\psi_0$ and
$\widehat M=U_B^\ast MU_B$, the normalised gradient~\eqref{eq:normgrad} is
\[
  \nabla\loss\;=\;\inner{\psi}{\,C\,\psi},
  \qquad C\;=\;[\widehat M,A]\;=\;\widehat MA-A\widehat M .
\]
The operator $C$ is selfadjoint, so $\nabla\loss\in\Qp$; and if
$M,A,U_A,U_B,\psi_0$ are integral then $\nabla\loss\in\Zp$, so
$\padicnorm{\nabla\loss}\le1$, equivalently
$\padicnorm{\partial_\theta\loss}\le p^{-1}$.
\end{lemma}

\begin{proof}
Since $U(\theta)\psi_0=U_B\exp_p(p\theta A)\psi$ and
$\inner{U_Bx}{MU_Bx}=\inner{x}{\widehat Mx}$, we have
$\loss(\theta)=\inner{\exp_p(p\theta A)\psi}{\widehat M\exp_p(p\theta A)\psi}$.
The coefficient of $\theta$ in the expansion is
\[
  \inner{pA\psi}{\widehat M\psi}+\inner{\psi}{\widehat M\,pA\psi}
  \;=\;p\,\inner{\psi}{(A^\ast\widehat M+\widehat MA)\psi},
\]
and $A^\ast=-A$ turns the bracket into $C=[\widehat M,A]$; dividing by $p$
gives the stated formula. For selfadjointness,
$C^\ast=(\widehat MA)^\ast-(A\widehat M)^\ast
=A^\ast\widehat M-\widehat MA^\ast=\widehat MA-A\widehat M=C$. For a
selfadjoint $C$ one has
$\overline{\inner\psi{C\psi}}=\inner{C\psi}{\psi}=\inner\psi{C^\ast\psi}
=\inner\psi{C\psi}$, so the scalar is fixed by the involution and lies in the
fixed field $\Qp$. Integrality is closure of $\OK$ under sums and products,
together with $\OK\cap\Qp=\Zp$.
\end{proof}

So the normalised gradient is a Hermitian form evaluated at a point of the
sphere. Both of its arguments are Haar-random: $\psi$ is uniform on $\Sph_N$ by
\Cref{lem:transitive}, and $C$ is determined by the block $U_B$. What remains
is a finite count, which we now do by hand.

\subsection{The counting lemma}

Fix $N\ge2$ and a selfadjoint $\Gamma$ over $\FF_{p^2}$. For $c\in\FF_p$ put
\[
  \mathcal{N}_c \;=\;
  \#\bigl\{u\in\FF_{p^2}^{N}\;:\;\inner uu=1,\ \inner{u}{\Gamma u}=c\bigr\}.
\]

\begin{definition}[Non-degeneracy]
\label{def:nondeg}
Say $\Gamma$ is \emph{$r$-non-degenerate} if it admits an orthonormal
eigenbasis over $\FF_{p^2}$ with all eigenvalues in $\FF_p$, and no eigenvalue
has multiplicity larger than $N-r$. Writing $m_{\max}$ for the largest
eigenvalue multiplicity, the second condition reads $r\le N-m_{\max}$.
\end{definition}

\begin{lemma}
\label{lem:count}
If $\Gamma$ is $r$-non-degenerate then for every $c\in\FF_p$
\[
  \left|\;\frac{\mathcal{N}_c}{|\widetilde\Sph_N|}\;-\;\frac1p\;\right|
  \;\le\;10\,p^{\,1-r},
  \qquad\text{and}\qquad
  \bigl|\widetilde\Sph_N\bigr| \;=\; p^{2N-1}-(-1)^{N}p^{N-1}.
\]
\end{lemma}

\begin{proof}
Work in an orthonormal eigenbasis of $\Gamma$, with eigenvalues
$\lambda_1,\dots,\lambda_N\in\FF_p$. The two constraints read
$\sum_i\Nm(u_i)=1$ and $\sum_i\lambda_i\Nm(u_i)=c$, where now
$\Nm:\FF_{p^2}\to\FF_p$, $\Nm(z)=z^{1+p}$. The fibres of $\Nm$ have sizes
$w(0)=1$ and $w(s)=p+1$ for $s\ne0$, since $\Nm$ restricts to a surjection
$\FF_{p^2}^\times\to\FF_p^\times$ with kernel of order $p+1$. Substituting
$t_i=\Nm(u_i)$,
\[
  \mathcal{N}_c\;=\;\sum_{\substack{t\in\FF_p^{N}\\ \sum t_i=1,\ \sum\lambda_it_i=c}}
  \ \prod_{i=1}^{N}w(t_i).
\]
Let $\chi$ be a non-trivial additive character of $\FF_p$ and detect the two
constraints by $\frac1p\sum_{\alpha}\chi(\alpha(\cdot))$. Then
\[
  \mathcal{N}_c\;=\;\frac1{p^2}\sum_{\alpha,\beta\in\FF_p}
  \chi(-\alpha-\beta c)\prod_{i=1}^{N}T(\alpha+\beta\lambda_i),
  \qquad
  T(s)=\sum_{t\in\FF_p}w(t)\chi(st).
\]
Direct evaluation gives $T(0)=1+(p+1)(p-1)=p^{2}$ and, for $s\ne0$,
$T(s)=1+(p+1)\sum_{t\ne0}\chi(st)=1-(p+1)=-p$. Hence, writing
$z(\alpha,\beta)=\#\{i:\alpha+\beta\lambda_i=0\}$,
\[
  \prod_iT(\alpha+\beta\lambda_i)\;=\;(p^{2})^{z}(-p)^{N-z}
  \;=\;(-1)^{N-z}p^{\,N+z}.
\]
Split the double sum into three ranges.

\textbf{Case $\alpha=\beta=0$.} Here $z=N$ and the term is $p^{2N}$.

\textbf{Case $\beta=0$, $\alpha\ne0$.} Here $z=0$, and
$\sum_{\alpha\ne0}\chi(-\alpha)=-1$; the total contribution is
$-(-1)^{N}p^{N}$.

\textbf{Case $\beta\ne0$.} For fixed $\beta\ne0$ the substitution
$\lambda=-\alpha/\beta$ is a bijection on $\alpha$. It gives
$\alpha+\beta\lambda_i=\beta(\lambda_i-\lambda)$, so $z=m(\lambda)$, the
multiplicity of $\lambda$ in the spectrum, zero if $\lambda$ is not an
eigenvalue; and $-\alpha-\beta c=\beta(\lambda-c)$. Summing over $\beta\ne0$
gives $\sum_{\beta\ne0}\chi(\beta(\lambda-c))=p\delta_{\lambda,c}-1$, so this
case contributes
$\Sigma(c)=\sum_{\lambda\in\FF_p}(-1)^{N-m(\lambda)}p^{\,N+m(\lambda)}
(p\delta_{\lambda,c}-1)$. Collecting,
\begin{equation}
\label{eq:collected}
  p^{2}\mathcal{N}_c\;=\;p^{2N}-(-1)^{N}p^{N}+\Sigma(c).
\end{equation}

We first extract the sphere count. Summing~\eqref{eq:collected} over $c\in\FF_p$
gives $\sum_c\mathcal{N}_c=|\widetilde\Sph_N|$ on the left, while on the right
$\sum_c\Sigma(c)=\sum_\lambda(-1)^{N-m(\lambda)}p^{N+m(\lambda)}(p-p)=0$;
therefore $p^{2}|\widetilde\Sph_N|=p^{2N+1}-(-1)^{N}p^{N+1}$, which is the
stated formula. In particular
$|\widetilde\Sph_N|\ge p^{2N-1}-p^{N-1}\ge\frac12p^{2N-1}$ for $N\ge2$, and
$|\widetilde\Sph_N|/p=p^{2N-2}-(-1)^{N}p^{N-2}$.

Now bound $\Sigma(c)$. At most $p$ values of $\lambda$ have $m(\lambda)=0$; for
those the summand has absolute value at most $p^{N}\cdot p=p^{N+1}$ when
$\lambda=c$ and at most $p^{N}$ otherwise, contributing at most
$p^{N+1}+(p-1)p^{N}\le2p^{N+1}$ in total. At most $p$ values of $\lambda$ have
$m(\lambda)\ge1$, and for those the summand is at most
$p^{\,N+m_{\max}}\cdot p$ in absolute value, contributing at most
$p^{\,N+m_{\max}+2}$. Hence, using $m_{\max}\ge1$,
\[
  \bigl|\,p^{2}\mathcal{N}_c-p^{2N}\,\bigr|
  \;\le\;p^{N}+2p^{N+1}+p^{\,N+m_{\max}+2}
  \;\le\;4\,p^{\,N+m_{\max}+2},
\]
that is $\mathcal{N}_c=p^{2N-2}+E$ with $|E|\le4p^{\,N+m_{\max}}$. Combining
with $|\widetilde\Sph_N|/p=p^{2N-2}-(-1)^{N}p^{N-2}$,
\[
  \left|\mathcal{N}_c-\frac{|\widetilde\Sph_N|}{p}\right|
  \;\le\;4p^{\,N+m_{\max}}+p^{\,N-2}\;\le\;5\,p^{\,N+m_{\max}},
\]
and dividing by $|\widetilde\Sph_N|\ge\frac12p^{2N-1}$ gives
$10\,p^{\,m_{\max}-N+1}\le10\,p^{\,1-r}$.
\end{proof}

\begin{remark}[When the bound bites]
\label{rem:bite}
\Cref{lem:count} is vacuous for small $r$: the deviation bound $10p^{1-r}$
exceeds $1$ unless $p^{\,r-1}>10$. Combined with the leading term $1/p$ in
\Cref{thm:main}, the conclusion is non-trivial as soon as
$p^{\,r-1}>10p/(p-1)$, which holds for $r\ge4$ when $p=3$, for $r\ge3$ when
$p\ge11$, and for $r\ge2$ when $p\ge13$. In the regime of interest $r$ grows
like $2^{n}$ (\Cref{ex:witness,rem:doubly}), so the condition is met from three
qubits onward for every odd $p$.
\end{remark}

\begin{remark}[The splitting clause is a real restriction]
\label{rem:notspectral}
It is tempting to read the first clause of \Cref{def:nondeg} as an analogue of
the complex spectral theorem, hence automatic. It is not. Over $\FF_{q^2}$ a
Hermitian matrix can have eigenvalues outside the fixed field $\FF_q$. Indeed
if $\Gamma v=\lambda v$ then
\[
  \lambda\inner vv=\inner v{\Gamma v}=\inner{\Gamma v}v=\bar\lambda\inner vv ,
\]
so $\lambda\ne\bar\lambda$ forces $\inner vv=0$; the same computation applied
to two vectors of the same eigenspace shows that eigenspace is totally
isotropic, and no orthonormal eigenbasis can exist. Concretely, take $q=3$ and
$\FF_9=\FF_3(\iota)$ with $\iota^2=-1$, so that the involution is
$x+\iota y\mapsto x-\iota y$ and $\Nm(x+\iota y)=x^2+y^2$. With $b=1+\iota$,
which has $\Nm(b)=2$, the matrix
\[
  \Gamma\;=\;\begin{pmatrix}0&b\\ \bar b&0\end{pmatrix}
  \;=\;\begin{pmatrix}0&1+\iota\\ 1-\iota&0\end{pmatrix}
\]
is Hermitian with characteristic polynomial $x^{2}-\Nm(b)=x^{2}+1$, whose roots
$\pm\iota$ lie in $\FF_9\setminus\FF_3$. Its $\iota$-eigenvector is
$(1+\iota,\iota)^{\mathsf T}$, of norm $2+1=0$.

So \Cref{def:nondeg} restricts to those $\Gamma$ whose characteristic
polynomial splits over $\FF_p$ with an orthonormal eigenbasis, and
\Cref{thm:main} is conditional on that. The general case needs the
classification of pencils of Hermitian forms over a finite field rather than a
single diagonalisation; we have not carried it out, and we do not know whether
the conclusion of \Cref{thm:main} survives without the hypothesis. This is the
principal gap in the paper. What we can say is that the restricted class is not
exotic: \Cref{ex:witness} exhibits a textbook observable generator pair inside
it, with $r$ as large as the pigeonhole bound of \Cref{rem:doubly} allows.
\end{remark}

\subsection{Main result}

\begin{theorem}[No barren plateaus]
\label{thm:main}
Let $p$ be odd, $K/\Qp$ unramified, $N=2^{n}$. Fix
\begin{itemize}
\item an integral observable $M=M^\ast$ and an integral generator
  $A=-A^\ast$, entering the ansatz~\eqref{eq:ansatz} as $\exp_p(p\theta A)$;
\item an admissible integral $\psi_0$;
\item an arbitrary $U_B\in G_N$, and write $\widehat M=U_B^\ast MU_B$,
  $C=[\widehat M,A]$ and $\Gamma=\widetilde C$;
\end{itemize}
and suppose $\Gamma$ is $r$-non-degenerate. Let $U_A$ be drawn from the Haar
measure $\nu$ on $G_N$. Then
\[
  \Pr\Bigl[\;\padicnorm{\partial_\theta\loss}=p^{-1}\;\Bigr]
  \;=\;\Pr\Bigl[\;\padicnorm{\nabla\loss}=1\;\Bigr]
  \;\ge\;1-\frac1p-10\,p^{\,1-r} .
\]
That is, the gradient attains the largest absolute value any convergent
$p$-adic ansatz permits (\Cref{rem:scale}), with a probability bounded below
independently of the number of qubits $n$.
\end{theorem}

\begin{proof}
By \Cref{lem:grad}, $\nabla\loss=\inner\psi{C\psi}\in\Zp$ with
$\psi=U_A\psi_0$, so $\padicnorm{\nabla\loss}\le1$ with equality precisely when
$v_p(\nabla\loss)=0$, that is when the residue
$\inner{\tilde\psi}{\Gamma\tilde\psi}\in\FF_p$ is non-zero.

The residue $\inner\psi{C\psi}\bmod p$ depends on $\psi$ only through
$\tilde\psi$, and $\Gamma$ is fixed by hypothesis. By \Cref{lem:transitive},
$\psi=U_A\psi_0$ is distributed according to the invariant measure on $\Sph_N$
as $U_A$ ranges over $G_N$, so $\tilde\psi$ is uniform on $\widetilde\Sph_N$.
Applying \Cref{lem:count} with $c=0$,
\[
  \Pr\bigl[\,v_p(\nabla\loss)\ge1\,\bigr]
  \;=\;\frac{\mathcal{N}_0}{|\widetilde\Sph_N|}
  \;\le\;\frac1p+10\,p^{\,1-r},
\]
and taking complements gives the claim. The identification of the two events in
the statement is $\nabla\loss=p^{-1}\partial_\theta\loss$.
\end{proof}

\begin{corollary}[Random second block]
\label{cor:randomUB}
If instead $U_A$ and $U_B$ are drawn independently from $\nu$ and $\Gamma$ is
$r$-non-degenerate for $\nu$-almost every $U_B$, the same bound holds.
\end{corollary}

\begin{proof}
The bound of \Cref{thm:main} is uniform over the admissible $U_B$, so it
survives averaging over $U_B$.
\end{proof}

Because \Cref{thm:main} takes a concrete pair $(M,A)$ as input, its hypothesis
is something one can check rather than assume. Here is a pair that passes.

% -----------------------------------------------------------------
% The following example is not strictly required by the argument; it exists
% to witness that the standing hypothesis of Thm 5.5 is satisfiable, with r
% exponentially large. Delete freely if space is short.
% -----------------------------------------------------------------
\begin{example}[A witness with $r=2^{\,n-1}$]
\label{ex:witness}
Let $p\equiv3\pmod4$, so that $\mu=-1$ is a non-square unit and
$K=\Qp(i)$ with $i^2=-1$, $\bar i=-i$. On $n$ qubits take the observable and
generator
\[
  M\;=\;X_1\;=\;\begin{pmatrix}0&\Id\\ \Id&0\end{pmatrix},
  \qquad
  A\;=\;-\tfrac{i}{2}\bigl(\Id-Z_1\bigr)
  \;=\;-i\begin{pmatrix}0&0\\ 0&\Id\end{pmatrix},
\]
blocks of size $2^{\,n-1}$; both are integral, since $2\in\Zp^\times$ for $p$
odd, and $M^\ast=M$ while $A^\ast=-A$ because $\bar i=-i$. In words: measure
Pauli $X$ on the first qubit and rotate the first qubit about $Z$. Take
$U_B=\Id$ and $\psi_0=e_1$. Then
\[
  C\;=\;[M,A]\;=\;\begin{pmatrix}0&-i\Id\\ i\Id&0\end{pmatrix}\;=\;Y_1 ,
\]
the $p$-adic Pauli $Y$ on the first qubit, and $\Gamma=\widetilde C$ satisfies
$\Gamma^2=\Id$. Its eigenvalues are therefore $\pm1\in\FF_p$, and since
$\Gamma$ has trace $0$ each occurs with multiplicity $2^{\,n-1}$. The
eigenspaces are $E_{\pm}=\{(x,\pm ix)^{\mathsf T}\}$ and carry the
non-degenerate forms $\inner{(x,\pm ix)}{(z,\pm iz)}=2\inner xz$, so each has
an orthonormal basis and $\Gamma$ is unitarily diagonalisable. Hence
$m_{\max}=2^{\,n-1}$ and $\Gamma$ is $r$-non-degenerate with
\[
  r\;=\;N-m_{\max}\;=\;2^{\,n-1},
\]
so \Cref{thm:main} gives $\Pr[\padicnorm{\partial_\theta\loss}=p^{-1}]\ge
1-p^{-1}-10\,p^{\,1-2^{n-1}}$. At $n=10$ and $p=3$ the correction term is
smaller than $3^{-500}$.
\end{example}

\begin{corollary}
\label{cor:deg}
Assume $\Gamma$ admits an orthonormal eigenbasis with eigenvalues in $\FF_p$.
Then $\Gamma$ fails to be $1$-non-degenerate if and only if $\widehat M$ and
$A$ commute modulo $p$.
\end{corollary}

\begin{proof}
Failure of $1$-non-degeneracy means $m_{\max}=N$, that is $\Gamma=\lambda\Id$
for some $\lambda\in\FF_p$. Being a commutator, $\Gamma$ has trace zero, so
$N\lambda=0$ in $\FF_p$; since $N=2^{n}$ and $p$ is odd, $p\nmid N$ and
therefore $\lambda=0$. Thus $\Gamma=0$, which is exactly
$[\widehat M,A]\equiv0\bmod\mm$. The converse is immediate.
\end{proof}

\begin{corollary}[Contrast with the complex case]
\label{cor:contrast}
For a complex $n$-qubit circuit whose blocks form $2$-designs,
$\Pr[\,|\partial_\theta\loss|\ge\varepsilon\,]\le O(2^{-n}\varepsilon^{-2})$ by
Chebyshev's inequality applied to the mean-zero, variance-$O(2^{-n})$ gradient
of~\cite{Mc18}. Over $K$, by \Cref{thm:main}, the gradient attains its
\emph{maximal} attainable absolute value with probability bounded below
uniformly in $n$.
\end{corollary}

\begin{remark}[How large can $r$ be, and how large is it typically]
\label{rem:doubly}
Two bounds frame the answer. Since $\Gamma$ has $N$ eigenvalues, all in the
$p$-element field $\FF_p$, pigeonhole gives $m_{\max}\ge\lceil N/p\rceil$ and
hence
\[
  r\;\le\;N-\lceil N/p\rceil\;\le\;N\Bigl(1-\frac1p\Bigr).
\]
A simple spectrum is thus impossible once $N>p$, which is the regime
$N=2^{n}$ reaches immediately; the correct picture is not ``$m_{\max}=O(1)$''
but ``the $N$ eigenvalues are spread over the $p$ available values''. If they
are spread evenly, $m_{\max}\approx N/p$ and $r\approx N(1-1/p)$. The witness
of \Cref{ex:witness} spreads them over two values and already achieves
$r=N/2$.

Either way $r=\Theta(N)=\Theta(2^{n})$, so the correction term $10p^{1-r}$ is
of order $p^{-2^{n}}$: doubly exponentially small in the qubit count. The
theorem is sharpest in exactly the regime -- many qubits, expressive circuits
-- where the complex barren plateau is most severe.
\end{remark}

% =================================================================
\section{Discussion}
\label{sec:discussion}

\subsection{What replaces the plateau}

The trainability obstruction does not vanish; it changes shape. In the complex
theory it is a \emph{measure} phenomenon: the bad set has full measure as
$n\to\infty$, and no amount of care in choosing $M$ and $A$ escapes it. Over
$K$, by \Cref{cor:deg}, the bad set is the locus where the observable and the
generator commute modulo $p$, cut out by the vanishing of a single matrix
commutator over $\FF_{p^2}$.

That locus is small but not null, and it is worth being precise about which.
By \Cref{lem:pushforward} the set $\{U_B:[\widehat M,A]\equiv0\bmod p\}$ is a
union of congruence classes, so it is open and closed in $G_N$ and its measure
is the proportion
\[
  \frac{\#\{g\in\Uni_N(\FF_{p^2}):[\,g^\ast\widetilde Mg,\widetilde A\,]=0\}}
       {|\Uni_N(\FF_{p^2})|},
\]
which is strictly positive whenever it is non-empty. It is a proper subvariety
of the level-one reduction, hence of measure $O(p^{-1})$ by a Schwartz--Zippel
count, but not of measure zero -- a distinction with no complex analogue, and
one that a finite-field ensemble makes unavoidable. The obstruction is
confined to a critical subvariety rather than spread across the landscape,
which matches the bifurcation structure Z\'u\~niga-Galindo obtains for very
deep classical $p$-adic networks~\cite{ZG26crit}, where the transition is
between a unique network state and infinitely many rather than between a flat
and a structured landscape.

There is also a design lesson, and it runs opposite to the complex one. In the
complex setting expressibility hurts: the more nearly the ansatz forms a
$2$-design, the flatter the landscape~\cite{Ho22}. \Cref{thm:main} inverts
this. Haar-randomness over $G_N$ is what \emph{produces} maximal gradients,
because a Haar-random residue $\inner{\tilde\psi}{\Gamma\tilde\psi}$ in $\FF_p$
is non-zero with probability close to $1-p^{-1}$; the danger is
under-randomised circuits whose commutators sit on the critical locus.
Choosing $\widehat M$ and $A$ so that they visibly fail to commute modulo $p$
is a concrete, checkable design criterion with no complex counterpart, and
\Cref{ex:witness} shows the most ordinary choice already satisfies it.

\subsection{Does a maximal gradient imply trainability?}

Not by itself. \Cref{thm:main} says the leading digit of the gradient is
generically non-zero, so a step $\theta\leftarrow\theta-\eta\nabla\loss$ with
$\padicnorm\eta=1$ moves the leading digit of the parameter. Convergence is
governed by a separate and stringent condition. $p$-adic gradient descent on a
quadratic loss with Hessian $H$ contracts at rate $\opnorm{\Id-\eta H}$, and
because the value group is discrete this is less than $1$ only when every entry
of $\Id-\eta H$ lies in $p\Zp$. For a \emph{scalar} step size that says
$H\equiv\eta^{-1}\Id\bmod p$: the Hessian must be a scalar matrix modulo $p$,
which is a far stronger demand than the real-analytic condition
$0<\eta<2/\lambda_{\max}$ it replaces. With a matrix preconditioner $P$ in
place of $\eta$ the requirement relaxes to invertibility of $H$ modulo $p$,
which is generic.

When contraction does hold, convergence is ultra-fast: the error contracts by
at least $p^{-1}$ per step, so $O(\log_p(1/\varepsilon))$ iterations suffice.
Our result removes gradient concentration from the list of obstructions. It
does not establish trainability, and the conditioning question is the natural
next one.

\subsection{Why a model with no hardware?}
\label{sec:why}

No $p$-adic quantum computer exists, and none is on any roadmap. It is worth
saying what a result about such a model is for.

The first use is diagnostic, and it concerns the complex theory rather than
this one. Barren plateaus are usually discussed as a phenomenon of quantum
circuits, attributed to entanglement, to depth, or to the expressibility of the
ansatz. The analysis above keeps the circuits, the observable, the ansatz shape
and the qubit count fixed, changes only the field the amplitudes live in, and
the phenomenon disappears. Whatever else barren plateaus are, they are not a
consequence of unitarity or of Hilbert-space dimension alone: the load is
carried by the compactness of the unitary group and by the Archimedean absolute
value, two properties of $\mathbb{C}$ that the physics never singles out. This
is a statement about the complex setting, and it suggests where a mitigation
has to bite. Restricting the ensemble so that it is not Haar-typical in
the full unitary group -- the mechanism behind structured ansatz design and
warm starts -- acts on an ingredient the argument genuinely needs. We do not
read this as a ranking of complex-side mitigations: the cost-function route
of~\cite{Ce21} attacks the same argument at a different point, and nothing here
says which attack is more effective.

The second use is methodological. Replacing $t$-design moment combinatorics by
character sums over finite unitary groups is not specific to the gradient.
\Cref{lem:pushforward} reduces \emph{any} Haar moment over $G_N$ to counting in
$\Uni_N(\OK/\mm^{k})$, so second-moment questions in this model -- concentration
of expectation values, the distribution of overlaps, the behaviour of a
$p$-adic kernel -- become finite computations of the same shape. The design
criterion the theorem yields has the same character: $[\widehat M,A]\not\equiv0
\bmod p$ is verified by reducing two integer matrices and multiplying them,
with no sampling and no asymptotics.

What the result does not do should be equally clear. It predicts nothing about
any existing device, and by \Cref{rem:nosampling} it cannot be checked by the
empirical protocol that produces barren-plateau plots in the complex setting;
the form of evidence that is available instead is set out under \emph{No
experiments} in \Cref{sec:limitations}. An application, if
one comes, would come through the classical side of the same programme, where
$p$-adic models are used because taxonomies, phylogenies and other hierarchical
data are ultrametric to begin with~\cite{Kh98,ZuGa22,ZG24dnn,Ma25}. The
contribution here is to that programme's theory, not to its hardware.

\subsection{Limitations}
\label{sec:limitations}

We state the boundaries of the result plainly.

\begin{itemize}
\item \textbf{Level one only.} \Cref{thm:main} controls
  $\Pr[v_p(\nabla\loss)\ge1]$. The full distribution of $v_p(\nabla\loss)$
  requires \Cref{lem:count} at level $\mm^{k}$ for $k>1$, over the finite rings
  $\OK/\mm^{k}$. \Cref{lem:pushforward} supplies the reduction; the character
  sum has not been carried out. We expect $\Pr[v_p\ge k]=p^{-k}(1+o(1))$, hence
  $\mathbb{E}[v_p(\nabla\loss)]=O(1)$, but we have not proved it.
\item \textbf{The splitting hypothesis.} This is the most serious limitation.
  By \Cref{rem:notspectral} the eigenbasis clause of \Cref{def:nondeg} is a
  genuine restriction and not an automatic consequence of Hermitian symmetry:
  over a finite field a Hermitian matrix may have eigenvalues outside the fixed
  field. \Cref{thm:main} therefore applies to those observable--generator pairs
  whose commutator splits, and we do not know what happens outside that class.
  \Cref{ex:witness} shows the class is non-empty and contains natural pairs;
  it does not show the class is large.
\item \textbf{Standing assumptions.} $p$ odd and $K/\Qp$ unramified. The
  ramified and $p=2$ cases change the residue field and the norm group and are
  not covered; $p=2$ is also the case in which the $p$-adic Hadamard gate does
  not exist~\cite{AMP23}.
\item \textbf{Single parameter, two blocks.} We differentiate one parameter in
  the two-block ansatz~\eqref{eq:ansatz}, following~\cite{Mc18}.
  Shallow-circuit and local-cost refinements in the manner of~\cite{Ce21} are
  open, as is the behaviour of the full parameter vector.
\item \textbf{No experiments.} By \Cref{rem:nosampling} the $p$-adic model
  admits no sampling-based readout, so the usual empirical protocol does not
  transfer. A meaningful numerical companion would be exact arithmetic in
  $\OK/\mm^{k}$ for small $n$ and $p$, verifying \Cref{lem:count} by
  enumeration; we regard this as the most useful immediate follow-up.
\end{itemize}

% =================================================================
\section{Related work}
\label{sec:related}

Barren plateaus were identified in~\cite{Mc18} and refined
in~\cite{Ce21,Ho22}; all of that analysis is Archimedean. The $p$-adic quantum
formalism we use is that of Aniello, Mancini and
Parisi~\cite{AMP23,AMP24,AMP23trace}, with the tensor product
of~\cite{Mu25tensor} and the qubit and gate constructions of Svampa and
co-authors~\cite{SMW22qub,SLMW26gates}; the exact algebraic sense of
universality goes back to Barenco et al.~\cite{Ba95}. Haar measure on $p$-adic
rotation groups was constructed in~\cite{AMP24haar} as an inverse limit of
measure spaces; our \Cref{lem:pushforward} is in the same spirit, and
\Cref{prop:noncompact} explains why such constructions must target compact
subgroups rather than the full unitary group, and hence why
\Cref{qst:plateau} is posed over $G_N$. On the learning side, classical $p$-adic neural networks
go back to~\cite{Kh98} and have been developed
extensively~\cite{ZuGa22,ZG24dnn,ZG26crit,Ma25}; the quantum bridge is so far
represented by the hierarchical quantum cellular network of~\cite{ZGZ26qnn}.
The generalised parameter-shift rules of~\cite{Cr19,Wi22} are the natural route
to a $p$-adic gradient estimator on the finite torsion group $\mu(K)$, which
our analytic $\exp_p$ formulation sidesteps but a hardware realisation would
need.

% =================================================================
\section{Conclusion}
\label{sec:conclusion}

Barren plateaus are a theorem about the Archimedean absolute value. Replacing
$\mathbb{C}$ by a quadratic extension of $\Qp$ removes them, at least on the
class of circuits whose observable--generator commutator splits over $\FF_p$.
There the gradient of a $p$-adic variational circuit generically attains the
largest absolute value the convergence radius allows, uniformly in the qubit
count, and the flat locus shrinks from a full-measure set to the subvariety
where the observable commutes with the generator modulo $p$. Whether the same
holds off that class is the first question we would ask next.

Getting there required two corrections and one substitution. The ensemble had
to be corrected: the full $p$-adic unitary group is non-compact and admits no
Haar probability measure, so the question posed over it has no answer.
The scale had to be corrected: every convergent $p$-adic ansatz confines its
one-parameter subgroups to the principal congruence subgroup, so $p^{-1}$, not
$1$, is what ``maximal gradient'' means. And the missing $t$-design
combinatorics had to be replaced by counting in finite unitary groups, a
substitution that appears reusable for other second-moment questions in
$p$-adic quantum machine learning. What remains is the conditioning of the
optimiser, which is now the binding constraint on trainability.

\bibliography{references}

\end{document}
